% Section: Setup (drafted 2026-10-30). Every claim cites its source in a comment; every number is \pv{...}.
\section{Setup}
\label{sec:setup}

\subsection{A generator of physics problems with known uncertainty}
\label{sec:generator}

% Spec (CLAUDE.md, component 1); physcal/env/graph.py; D51 (generator 2.0.0); configs/env_v2.yaml.
The benchmark is a procedural generator of multi-step physics word problems (generator version \pv{gen.version}).
It has \pv{gen.nfamilies} formula families (\pv{gen.families}), each one physical scenario with named, dimensioned
quantities and a set of relations written in SymPy~\citep{meurer2017sympy}: \pv{gen.nformulas.kinematics},
\pv{gen.nformulas.dynamics} and \pv{gen.nformulas.energy} relations respectively.
% physcal/env/graph.py: Family.sample_world, chains, is_valid_chain.
A problem is a linear chain through its family's formula graph.
The generator samples a consistent world (basis quantities log-uniform within physical ranges, derived quantities
checked against plausibility bounds), picks a chain, shows the chain's inputs (the \emph{givens}) and asks for its
output (the \emph{target}).
The depth of a chain, from one step up to \pv{gen.maxdepth}, is its difficulty dial.
% graph.py docstring: "the target's minimal derivation depth ... must equal the chain length ... no given may be
% derivable from the other givens".
Depth is honest: the target's shortest derivation from the givens, over all relations of the family, has exactly
the chain's length, and no given can be derived from the others.
% env_v2 splits.heldout_depths; R5-8 (depth 5 is extrapolation in chain length, not to new physics).
Depth \pv{gen.heldoutdepth} is never used for training.
It tests extrapolation in chain length within the same families, not transfer to new physical laws.

% D10 (distractors), D11 (plausibility bounds), CLAUDE.md component 1 (units, phrasing), env_v2 display settings.
Each problem's text is randomized: phrasing templates, display units with SI prefixes and non-SI units, and
\pv{gen.distractors} distractor quantities that the answer does not depend on.
% env_v2 max_condition_number; physcal/env/graph.py condition_number.
Chains whose output amplifies relative errors in a given by more than a factor of \pv{gen.kappa}
(the largest $|\partial \ln(\text{output}) / \partial \ln(\text{given})|$) are rejected.
% physcal/env/graph.py evaluate_chain_with_units (pint dimension checks).
The ground truth is computed in SI units with dimension checks at every step, then converted to the requested unit.
Every problem is reproducible from its generator version and seed.

\subsection{Measurement noise and the true uncertainty $\sigstar$}
\label{sec:sigstar}

% D51 (noise mode: 5 s.f. so that rounding is negligible against the noise, FINDINGS #8; a problem without noise
% behaves as generator 1.1.0, 3 s.f., truth from the displayed values); D113 (mixed levels: every given at 5 s.f.,
% level-0 givens exact); D54 (levels capped at 5% after FINDINGS #9; a post-hoc design change); configs/pairs.yaml.
In the measurement-noise mode, each noisy given is shown as a measurement
$x_{\text{meas}} = x_{\text{true}}(1+e)$, $e \sim \mathcal{N}(0, s^2)$, with \pv{gen.noisesigfigs} significant
figures, so that rounding is negligible against the noise, and a stated relative uncertainty ``$\pm\, s$''; a value
shown without one is exact.
The truth is the noise-free chain evaluated on the unrounded true values.
A problem without noise shows \pv{gen.sigfigs} significant figures, and its truth is computed from the displayed
values.
The pair sets of Section~\ref{sec:pairs} use the levels \pv{gen.levels}, and the training data of
Section~\ref{sec:dpos} draws a level for each given independently from \pv{gen.mixedlevels}, where a level of zero
means an exact given.

% D50/A2 (sigma* definition), D51 (stored per problem; n_mc draws); JCGM 101 (propagation of distributions).
Because the generator knows the measurement model, it knows each problem's correct uncertainty.
Given the displayed values, the clean value of each noisy given is drawn as $x = x_{\text{meas}}/(1+e')$,
$e' \sim \mathcal{N}(0, s^2)$, with exact givens and constants held fixed; \pv{gen.nmc} such draws are pushed through
the chain, as in Monte Carlo propagation of distributions~\citep{jcgm101}.
The target $\sigstar$ is the scale of the Gaussian in log space, centred at $\ln f(x_{\text{meas}})$, that minimizes
the expected CRPS against these draws:
\begin{equation}
  \sigstar = \arg\min_{\sigma > 0}\; \mathbb{E}_{x}\!\left[\mathrm{CRPS}\!\left(\mathcal{N}(\ln f(x_{\text{meas}}),
  \sigma^2),\, \ln f(x)\right)\right],
  \label{eq:sigstar}
\end{equation}
where $f$ is the chain and $f(x_{\text{meas}})$ is what a correct solver computes from the displayed values.
% D50/A2: also stored, the SD of the log draws and the delta-method value (used to build the pair sets).
The standard deviation of the log draws and the first-order (delta-method) value are stored as well.

% D50/A3 (gate), D57 (miss-count rule; FINAL run), FINDINGS #12 (passed), #9 and #11 (the two failed runs);
% results/sigma_star_gate_final/gate.json.
Before any model was evaluated on noisy problems, the generator had to pass a self-consistency gate: intervals
$\ln f(x_{\text{meas}}) \pm k\sigstar$ must cover the clean truth at their nominal rates.
On the final run, with \pv{gate.n} noisy problems in \pv{gate.ncells} depth-by-level cells of \pv{gate.percell}
each, pooled coverage was \pv{gate.cov.one} at $k=1$ (nominal \pv{gate.nominal.one}, Kupiec~\citep{kupiec1995}
$p = \pv{gate.kupiec.one}$) and \pv{gate.cov.two} at $k=2$ (nominal \pv{gate.nominal.two},
$p = \pv{gate.kupiec.two}$).
No cell was rejected after Holm correction, and the probability integral transform of the truth was uniform
(Kolmogorov--Smirnov $p = \pv{gate.pitp}$).
The \pv{gate.nzero} noise-free control problems were exact.
% FINDINGS #9 (failed on one cell at 10% noise) -> D54 (cap at 5%, after the result); FINDINGS #11 (failed only on
% the cell-share rule, which demanded zero misses) -> D57 (miss-count test, before the final run).
The gate failed twice before this run.
The first run failed on one cell at the highest noise level, after which the noise levels were capped; that cap was
chosen after seeing the result.
The second failed only on a rule that, as written, tolerated no cell outside its interval; it was replaced by a
miss-count test before the final run.
Both changes are listed in Table~\ref{tab:deviations}.

\subsection{Output format and verifier}
\label{sec:verifier}

% CLAUDE.md component 3 (format); D16, D17 (parser); D83 (the <confidence> channel, verifier 1.2.0).
A completion may reason freely, then reports the requested intermediate values in a \texttt{<steps>} block and the
answer as \texttt{<answer>}$\mu \pm \sigma$ \textit{unit}\texttt{</answer>}.
% D13 (failures score the rule's floor); CLAUDE.md component 4 (units gate).
A completion that does not parse, or whose unit has the wrong dimension, scores the rule's floor
(\pv{ver.floor}); an answer in another unit of the right dimension is converted.

% D15 (the spec's relative CRPS is only approximately proper), D29 (crps_log is the default), D38 (its sigma floor);
% configs/reward.yaml; Gneiting & Raftery (2007); closed form of the Gaussian CRPS: Gneiting et al. (2005).
The scoring rule, $\crpslog$, scores the log of the truth $y$ against $\mathcal{N}(\ln\mu, (\sigma/\mu)^2)$:
\begin{equation}
  \crpslog(\mu,\sigma; y) = -\,s\left[z\left(2\Phi(z)-1\right) + 2\varphi(z) - \tfrac{1}{\sqrt{\pi}}\right],\qquad
  z = \frac{\ln y - \ln \mu}{s},\quad s = \max\!\left(\sigma/\mu,\ \epsilon\right),
  \label{eq:crpslog}
\end{equation}
the Gaussian CRPS in closed form~\citep{gneiting2005emos}, clipped below at the floor.
It is strictly proper for log-normal beliefs~\citep{gneiting2007strictly} and scale-free.
% D15: dividing by the truth reweights outcomes by 1/|y|; the optimal mu shifts and the optimal sigma inflates.
The relative CRPS of the original specification divides by the truth, which reweights outcomes and makes it only
approximately proper; it remains an ablation option.
% D37, D38: a floor of 0.1% made the reward flat in sigma for an overconfident policy; verifier 1.1.0 floors at 1e-6.
The floor $\epsilon$ is \pv{ver.sigmafloor} (verifier \pv{ver.version}): an earlier, larger floor made the reward
flat in $\sigma$ exactly where an overconfident policy needed a gradient.
% configs/reward.yaml weights and steps; D14 (step credit can never outweigh the final answer).
During reinforcement learning the reward adds \pv{ver.stepweight} times the share of requested intermediates within
\pv{ver.steptol} of their values; this weight is checked to be too small to ever outweigh the final score.
An answer is \emph{accurate} if it is within \pv{ver.binary} of the truth.

% D50/A1 (computation-correct subset); configs/pairs.yaml analysis.tau.
On noisy problems, an answer is \emph{computation-correct} if $|\ln\mu - \ln f(x_{\text{meas}})| \leq
\pv{ver.tau}$, that is, if it is what a correct solver computes from the displayed values.
Comparisons of a stated $\sigma$ with $\sigstar$ that require a correct answer use this subset.

\subsection{The oracle-centred expected score}
\label{sec:expected}

% FINDINGS #24 (stage A), D100; results/stage_a/results.json; results/r5_diag/tail.json (the gap).
A single realized truth says little about whether a stated $\sigma$ is right: a $\sigma$ that is \pv{sa.scale}
times too small scores better than $\sigstar$ against the realized truth on \pv{sa.halfwins} of problems.
Because the generator knows the predictive distribution, we also score a stated relative $\sigma$ by its
\emph{expected} $\crpslog$ over \pv{sa.draws} fresh predictive draws, centred at $f(x_{\text{meas}})$.
Per-problem differences between sources then have \pv{sa.varratio.min}--\pv{sa.varratio.max} times less variance than
the realized ones.
Expected-score differences are reported with their intervals and, where useful, as a share of the gap between the
noise heuristic RSS (Section~\ref{sec:comparators}) and $\sigstar$, which is \pv{gap.rss} on the frozen
confirmatory set.

\subsection{Counterfactual pair sets}
\label{sec:pairs}

% D50/A6, D58 (A6), D59 (A6b), D95 (A6c), D113 (the frozen confirmatory set), D147 (the J set);
% configs/pairs.yaml; Table tab:pairs is generated from the sets' manifests.
A model can state a plausible $\sigma$ without propagating anything, by learning which templates, depths and stated
percentages go with which uncertainty.
To separate propagation from such a lookup, we evaluate on matched pairs of problems that differ in exactly one
respect (Table~\ref{tab:pairs}).
The statistic is the within-pair slope: the ordinary least-squares slope, with an intercept, of
$\Delta\ln\sigma_{\text{stated}}$ on $\Delta\ln\sigstar$ across pairs, where $\Delta$ is the difference between a
pair's two members.
An exact propagator has slope one; a reader of features that do not differ within a pair has slope zero.

\begin{description}[leftmargin=0pt, itemsep=2pt]
% D58; pairs.yaml construction.
\item[A6 (value changes).] One chain template, two worlds: of \pv{pairs.aSix.worlds} candidate worlds, the two with
  the lowest and highest delta-method $\sigstar$, kept if their ratio is at least \pv{pairs.aSix.ratio}.
  Template, phrasing, noise level and the standardized measurement errors are shared.
  Power-law chains cannot form A6 pairs, since their relative $\sigstar$ does not depend on the values.
% D59; pairs.yaml invariance_construction.
\item[A6b (invariance).] Power-law templates only, two worlds whose targets differ by a factor of at least
  \pv{pairs.aSixB.ratio}, with relative $\sigstar$ checked equal by Monte Carlo (to within \pv{pairs.aSixB.mc} in
  log).
  Here the slope is taken on $\Delta\ln y$ and must be zero: a model whose relative $\sigma$ grows with the size of
  the answer fails it.
% D95; pairs.yaml allocation_construction; FINDINGS #23 (no surface reader propagates on A6c).
\item[A6c (noise allocation).] One chain, one world and identical displayed values.
  The two members differ only in which given carries the higher of two noise levels: in one member it sits where it
  moves $\sigstar$ least, in the other where it moves it most (ratio at least \pv{pairs.aSixC.ratio}).
  Both members show the same multiset of percentages, so no reader of the problem's surface can tell them apart.
\end{description}

% pairs.yaml export (cells x pairs); D113 sec. 8 (frozen confirmatory set; loader guard).
Each development set has \pv{pairs.cells} depth-by-level cells of \pv{pairs.percell} pairs.
The frozen confirmatory set (Section~\ref{sec:frozen}) repeats the three constructions with fresh seeds,
\pv{frozen.aSixC.problems} problems each.

% D147 (pre-registration), D149 (build); data/jacobian_support/r4_1.json; confirmatory/r5_6 manifests.
\paragraph{Jacobians outside the training range (the J set).}
A model trained on a chain might have learned that chain's typical sensitivities rather than computing them.
The J set therefore uses only chains that occur in the training data, and moves them to Jacobians they never had
there.
For each of the \pv{j.support.chains} chains in the \pv{rfour.n} training problems we record the range of every
given's log-sensitivity $\partial\ln f/\partial\ln x$ at the displayed values.
A J world has at least one given whose sensitivity lies outside that range by at least \pv{j.margin}, while every
displayed and computed value stays inside the range that quantity had in training; the first such world of up to
\pv{j.worlds} draws is used.
\emph{J sensitivity pairs} apply the A6c construction to J worlds.
\emph{J invariance pairs} pair a J world with a copy rescaled along a dimensional direction that leaves every
physical constant unchanged, which keeps every log-sensitivity, and so the relative $\sigstar$, identical while the
target changes by a factor of \pv{j.scale.min} to \pv{j.scale.max}; A6b cannot be used here because its power-law
Jacobians are constant.
% J manifests: median exceedance and median chain training count; D147 sec. 10.
The J set has \pv{j.sens.problems} sensitivity and \pv{j.inv.problems} invariance problems at the trained depths
(up to \pv{j.maxdepth}). Across the sensitivity pairs, the most out-of-range sensitivity lies a median of
\pv{j.sens.exc} outside its chain's range, and the median chain was seen \pv{j.sens.ntrain} times in training.

% GENERATED by scripts/paper/tables.py; do not edit. Rebuild: uv run python scripts/paper/build.py
% artifact: table:pair_sets
% sources (path @ git blob):
%   configs/pairs.yaml @ d305fb030896
%   data/eval_sets/2.0.0+bdbb2279/a6_pairs_manifest.json @ 297b24dc5d09
%   data/eval_sets/2.0.0+bdbb2279/a6b_pairs_manifest.json @ bb40ea9f03dc
%   data/eval_sets/2.0.0+bdbb2279/a6c_pairs_manifest.json @ bb2ebd3748b8
%   data/eval_sets/2.0.0+bdbb2279/confirmatory/a6_pairs_manifest.json @ a4667a2fef54
%   data/eval_sets/2.0.0+bdbb2279/confirmatory/a6b_pairs_manifest.json @ f42552284794
%   data/eval_sets/2.0.0+bdbb2279/confirmatory/a6c_pairs_manifest.json @ e413b9b471bf
%   data/eval_sets/2.0.0+bdbb2279/confirmatory/r5_6/a6b_pairs_manifest.json @ 33486465e0b3
%   data/eval_sets/2.0.0+bdbb2279/confirmatory/r5_6/a6c_pairs_manifest.json @ 3e36aa054364
\begin{table}[t]
\centering
\small
\caption{The counterfactual pair sets and what each isolates.}
\label{tab:pairs}
\begin{tabular}{@{}L{2.1cm}L{3.6cm}L{3.1cm}cc r@{}}
\toprule
set & changes within a pair & held fixed & regressor & ideal & problems \\
\midrule
A6 & the world (values); same template & template, phrasing, level, $z$ & $\Delta\ln\sigma^\star$ & 1 & \num{1020} / \num{1020} \\
A6b & the target's magnitude ($\geq\times$\num{2}); power-law template & relative $\sigma^\star$ (MC-checked) & $\Delta\ln y$ & 0 & \num{1020} / \num{1020} \\
A6c & which given carries the high level & displayed values, level multiset & $\Delta\ln\sigma^\star$ & 1 & \num{1020} / \num{1020} \\
J sensitivity & as A6c, on a Jacobian outside the training range & as A6c & $\Delta\ln\sigma^\star$ & 1 & \num{816} \\
J invariance & a dimensional rescaling (target $\times$\num{2}--\num{8}) & every log-sensitivity, so relative $\sigma^\star$ & $\Delta\ln y$ & 0 & \num{816} \\
\bottomrule
\end{tabular}
\par\smallskip\begin{minipage}{0.95\linewidth}\footnotesize Ideal: the within-pair OLS slope of $\Delta\ln$(stated $\sigma$) on the regressor for an exact propagator; a reader of surface features that are identical within a pair scores 0 on A6c and J. Problems: development / frozen confirmatory set (A6, A6b, A6c; D113), or the frozen J set (D147). Sources: the sets' manifests.\end{minipage}
\end{table}

